\section*{پیش از شروع}

\begin{blockbox}{قواعد، یادداشت‌ها و سیاست‌ها}
    \begin{enumerate}
        \item پیش از شروع، دست‌کم یک بار اسلایدهای درس را مرور کنید.

        \item استفاده از ابزارهای هوش مصنوعی برای جست‌وجو و مطالعه آزاد است.
              سپردن نگارش پاسخ نهایی به آن‌ها آزاد نیست و هر سؤالی که این‌گونه
              نوشته شده باشد نمره‌ی صفر می‌گیرد.

        \item گفت‌وگو درباره‌ی سؤال‌ها با هم‌کلاسی‌ها تشویق می‌شود؛ رد و بدل
              کردن پاسخ نه، و طبق آیین‌نامه‌ی امانت‌داری علمی دانشگاه برای هر
              دو طرف رسیدگی می‌شود.

        \item این تمرین \hwgrade{} نمره از نمره‌ی درس دارد و برای دقت بیشتر از
              ۱۰۰ تصحیح می‌شود.

        \item هر کتاب، مقاله و مخزنی را که از آن استفاده کرده‌اید در انتهای
              پاسخ‌ها ذکر کنید.

        \item سیاست تأخیر همان است که ابتدای نیم‌سال توافق شد.
    \end{enumerate}
\end{blockbox}

\subsection*{چه چیزی تحویل می‌دهید}

یک آرشیو \lr{\texttt{.zip}} با این نام و ساختار:

\begin{codeblock}{submission.txt}
@\coursecode@-HW@\hwnumber@-[STDID].zip
|
|-- @\coursecode@-HW@\hwnumber@-[STDID].pdf
|-- HW@\hwnumber@-Practical/
        |-- *
\end{codeblock}

که \lr{\texttt{STDID}} شماره‌ی دانشجویی شماست. فایل \lr{PDF} شامل پاسخ‌های
نظری به همراه گزارش بخش عملی است و فایل‌های عملی در پوشه می‌روند.

\subsection*{پیش از بارگذاری بررسی کنید}

\begin{checklist}
    \citem به همه‌ی سؤال‌ها پاسخ داده‌اید، یا صراحتاً نوشته‌اید که پاسخ نداده‌اید.
    \item گزارش، فهرست کتابخانه‌های استفاده‌شده و دلیل هرکدام را دارد.
    \item کد از روی یک نسخه‌ی تازه‌ی مخزن اجرا می‌شود.
    \item نام آرشیو دقیقاً مطابق قالب بالاست.
\end{checklist}

\begin{asidebox}{درباره‌ی نمره‌ی جزئی}
    پاسخ ناقصِ صادقانه‌ای که مسیر حل را نشان دهد، بیش از پاسخ قاطعِ نادرست
    نمره می‌گیرد. بنویسید چه چیزی را امتحان کرده‌اید و کجا شکست خورده است.
\end{asidebox}

موفق باشید.
